\subsection{极分解}

在本节中，我们考虑复 Hilbert 空间。

设 \(E_{1}\) 与 \(E_{2}\) 为 Hilbert 空间，\(u \in \mathcal{L}(E_{1}; E_{2})\) 。自同态 \(u^{*}u\) （\(E_{1}\) 的自同态）是正的（命题 8），因而可以作正元素 \((u^{*}u)^{1/2}\) ，它属于 \(\mathcal{L}(E_{1})\) 。

定义 3.　我们称 \((u^{*}u)^{1/2}\) 为 u 的绝对值，并把它记作 \(|u|\) 。

在 \(E_{1} = E_{2}\) 的情形，这个定义与 I，第 119 页，注记 9 中给出的定义一致。

对元素 \(u\) （\(\mathcal{L}(\mathrm{E}_{1};\mathrm{E}_{2})\) 的元素），回忆（EVT，V，第 41 页，定义 2）：\(u\) 的初始子空间是闭子空间 \(\mathrm{Ker}(u)^{\circ}\) （\(\mathrm{E}_{1}\) 中的闭子空间），\(u\) 的最终子空间是闭子空间 \(\overline{\mathrm{Im}(u)}\) （\(\mathrm{E}_{2}\) 中的闭子空间）。

命题 10.　设 \(\mathrm{E}_1\) 与 \(\mathrm{E}_2\) 为复 Hilbert 空间，\(u \in \mathcal{L}(\mathrm{E}_1; \mathrm{E}_2)\) 。

\begin{enumerate}
\def\labelenumi{\alph{enumi})}
\item
  \(|u|\) 的初始子空间与最终子空间都等于 \(u\) 的初始子空间，且我们有 \(\| |u| \| = \| u\|\) ；
\item
  存在唯一的部分等距 \(j\) ，它从 \(\mathrm{E}_1\) 到 \(\mathrm{E}_2\) 中，满足 \(\operatorname{Ker}(j) = \operatorname{Ker}(u)\) 与 \(u = j|u|\) ；
\item
  j 的初始（相应地，最终）子空间等于 u 的初始（相应地，最终）子空间；
\item
  设 \(u_{1}\) 是 \(\mathcal{L}(\mathrm{E}_{1})\) 的正元素，\(j_{1}\) 是 \(\mathcal{L}(\mathrm{E}_{1};\mathrm{E}_{2})\) 的部分等距，满足 \(u = j_{1}u_{1}\) 与 \(\operatorname{Ker}(j_{1}) = \operatorname{Ker}(u_{1})\) 。则 \(u_{1} = |u|\) 且 \(j_{1} = j\) 。
\end{enumerate}

对每个 \(x\in \mathrm{E}_1\) ，有

\[
\| u (x) \| ^ {2} = \langle x | (u ^ {*} u) (x) \rangle = \langle x | | u | ^ {2} (x) \rangle = \| | u | (x) \| ^ {2}.\tag{2}
\]

这就证明 \(\operatorname{Ker}(u) = \operatorname{Ker}(|u|)\) 与 \(\| |u| \| = \|u\|\) 。由于 \(|u|\) 是 Hermite 的，\(|u|\) 的像的闭包是其核的正交补（EVT，V，第 41 页，命题 4），即 \(u\) 的初始空间，由此得 \(a\) ）。

公式 (2) 蕴涵：存在等距映射 \(v\) ，它从 \(\operatorname{Im}(|u|)\) 到 \(\operatorname{Im}(u)\) 上，使得 \(v(|u|(x)) = u(x)\) 对所有 \(x \in \mathrm{E}_1\) 成立。设 \(j\) 为 \(\mathcal{L}(\mathrm{E}_1; \mathrm{E}_2)\) 中唯一扩张 \(v\) 且在 \(\operatorname{Im}(|u|)^{\circ} = \operatorname{Ker}(|u|) = \operatorname{Ker}(u)\) 上为零的元素。则 \(j\) 具有 \(b\) ）中的性质。\(j\) 的唯一性由分解 \(\mathrm{E} = \operatorname{Ker}(u) \oplus \overline{\operatorname{Im}(|u|)}\) 得出。

\(j\) 的初始子空间是 \(\mathrm{Ker}(j)^{\circ} = \mathrm{Ker}(u)^{\circ}\) ，即 \(u\) 的初始空间。它的最终子空间是 \(\overline{j(\mathrm{Ker}(u)^{\circ})} = \overline{j(\mathrm{Im}(|u|))} = \overline{\mathrm{Im}(u)}\) ，即 \(u\) 的最终子空间。这就证明了 \(c\) ）。

现在设 \(u_1\) 与 \(j_1\) 如 \(d\) ）中所述。我们有 \(u^*u = u_1j_1^*j_1u_1\) 。映射 \(j_1^*j_1\) 是以 \(\text{Ker}(j_1) = \text{Ker}(u_1)\) 为核的正交投影（EVT，V，第 41 页，命题 5 (ii)），因而以 \(\overline{\text{Im}(u_1)}\) 为像。因此 \(u^*u = u_1^2\) ，从而 \(u_1 = (u^*u)^{1/2} = |u|\) （I，第 118 页，命题 16）。最后，\(b\) ）中的唯一性断言蕴涵 \(j_1 = j\) 。

定义 4.　设 \(E_{1}\) 与 \(E_{2}\) 为复 Hilbert 空间，\(u \in \mathcal{L}(E_{1}; E_{2})\) 。我们称偶 \((j, |u|)\) 为 u 的极分解，其中 j 是从 \(E_{1}\) 到 \(E_{2}\) 中的唯一的部分等距映射，满足 \(u = j|u|\) 与 \(\operatorname{Ker}(j) = \operatorname{Ker}(u)\) 。

命题 11.　设 \(\mathrm{E}_1\) 与 \(\mathrm{E}_2\) 为复 Hilbert 空间，\(u \in \mathcal{L}(\mathrm{E}_1; \mathrm{E}_2)\) 。设 \((j, |u|)\) 为 \(u\) 的极分解。

\begin{enumerate}
\def\labelenumi{\alph{enumi})}
\item
  我们有 \(|u| = j^{*}u = u^{*}j\) ；
\item
  我们有 \(|u^{*}| = j u^{*} = u j^{*}\) ；
\item
  \(u^{*}\) 的极分解是 \((j^{*}, |u^{*}|)\) 。
\end{enumerate}

设 I = \(\operatorname{Ker}(u)^{\circ}\) 与 \(F = \overline{\operatorname{Im}(u)}\) 为 u 的初始子空间与最终子空间；此外，我们有 I = \(\operatorname{Ker}(|u|)^{\circ} = \overline{\operatorname{Im}(|u|)}\) （命题 10，a)）。映射 \(j^{*}j\) 是 \(E_{1}\) 到 I 上的正交投影（同前及

EVT，V，第 41 页，命题 5 (ii)）。因此我们有 \(j^{*}u = j^{*}j|u| = |u|\) ，进而 \(u^{*}j = (j^{*}u)^{*} = |u|^{*} = |u|\) ，由此得 \(a\) ）。

类似地，我们算得 \(u^{*} = |u|j^{*} = (j^{*}j|u|)j^{*} = j^{*}(j|u|j^{*})\) 。自同态 \(j|u|j^{*}\) （\(E_{2}\) 的自同态）是正的。线性映射 \(j^{*}\) 是部分等距的，其初始空间为 F，最终空间为 I（EVT，V，第 41 页，命题 5），且由 j 与 \(j^{*}\) 通过过渡到子空间而得到的从 I 到 F（相应地，从 F 到 I）的线性映射是互逆的同构（同前）。因此我们有 \(\operatorname{Ker}(j|u|j^{*}) = \operatorname{Ker}(|u|j^{*}) = \operatorname{Ker}(j^{*})\) ，因为 \(j^{*}\) 的像包含于 \(\operatorname{Ker}(u)^{\circ} = \operatorname{Ker}(|u|)^{\circ}\) 。由命题 10，d)，偶 \((j^{*}, j|u|j^{*})\) 是 \(u^{*}\) 的极分解。这就证明了 c)，而断言 b) 则由把断言 a) 应用于 \(u^{*}\) 得出。

推论.　设 \(E_{1}\) 与 \(E_{2}\) 为复 Hilbert 空间，\(u \in \mathcal{L}(E_{1}; E_{2})\) 。我们有 \(\operatorname{Im}(u) = \operatorname{Im}(|u^{*}|)\) 。

设 \((j, |u|)\) 为 \(u\) 的极分解。由前一个命题，我们有 \(|u^{*}| = j|u|j^{*} = uj^{*}\) 。由 EVT，V，第 41 页，命题 5，映射 \(j^{*}\) 是部分等距的。其最终空间为 \(\operatorname{Ker}(j)^{\circ} = \operatorname{Ker}(u)^{\circ}\) （命题 10，c)）。断言得证。

命题 12.　设 \(E_{1}\) 与 \(E_{2}\) 为复 Hilbert 空间，\(u \in \mathcal{L}(E_{1}; E_{2})\) 。设 \((j, |u|)\) 为 u 的极分解。u 是双射的充分必要条件是：\(|u|\) 在 \(\mathcal{L}(E_{1})\) 中可逆，且 j 是从 \(E_{1}\) 到 \(E_{2}\) 上的同构。

这个条件是充分的。反之，若 u 是双射，则 \(u^{*}u\) 在 \(\mathcal{L}(\mathrm{E}_{1})\) 中可逆，且 \(|u| = (u^{*}u)^{1/2}\) 也是双射，因为它的谱包含于 \(R_{+}^{*}\) 。此外，\(\operatorname{Ker}(j) = \operatorname{Ker}(u) = \{0\}\) 且 \(\operatorname{Im}(j) = \operatorname{Im}(u) = F\) ，故 j 把 \(E_{1}\) 等距地映到 \(E_{2}\) 上。

命题 13.　设 E 为复 Hilbert 空间，u 为 E 的自同态。下列条件等价：

\begin{enumerate}
\def\labelenumi{(\roman{enumi})}
\item
  自同态 u 是正规的；
\item
  存在 \(\mathcal{L}(\mathrm{E})\) 的酉元素 v，它与 \(|u|\) 交换，使得 \(u = v|u|\) 。
\end{enumerate}

设 \((j,|u|)\) 为 \(u\) 的极分解。假设 \(u\) 是正规的。这时我们有 \(|u^{*}| = (uu^{*})^{1/2} = (u^{*}u)^{1/2} = |u|\) 。于是命题 11 蕴涵 \(|u|j = |u^{*}|j = ju^{*}j = j|u|\) 。此外，\(j\) 使互补的正交子空间 \(\operatorname{Ker}(|u|) = \operatorname{Ker}(j)\) 与 \(\overline{\operatorname{Im}(|u|)} = \overline{\operatorname{Im}(j)}\) 保持稳定（命题 10）。设 \(v\) 为 \(\mathcal{L}(\mathrm{E})\) 中与 \(j\) 在 \(\operatorname{Ker}(u)^{\circ}\) 上、与恒等映射在 \(\operatorname{Ker}(u)\) 上相一致的元素。由于 \(j\) 诱导出从 \(\operatorname{Ker}(u)^{\circ}\) 到 \(\overline{\operatorname{Im}(u)} = \operatorname{Ker}(u^{*})^{\circ} = \operatorname{Ker}(u)^{\circ}\) 上的等距（因为 \(u\) 是正规的），自同态 \(v\) 是酉的；此外，它与 \(|u|\) 交换，因为 \(j|u| = |u|j\) ，且我们有 \(u = v|u|\) 。

反之，设 v 是 \(\mathcal{L}(\mathrm{E})\) 的酉元素，它与 \(|u|\) 交换，使得 \(u = v|u|\) 。我们有 \(uu^{*} = v|u|^{2}v^{*} = |u|^{2}vv^{*} = |u|^{2} = u^{*}u\) ，因此 u 是正规的。

\pandocbounded{\includegraphics[keepaspectratio]{images/part001_8ac1eb40.jpg}}

设 E 为复 Hilbert 空间，\(u \in \mathcal{L}(\mathrm{E})\) 。设 \((j, |u|)\) 为 u 的极分解。j 与 \(|u|\) 交换而 u 不是正规的情形是可能的（I，第 189 页，习题 11）。

